\documentclass{article}
\usepackage{amsmath,url}
\title{Can a charged probe identify which receiver an outstanding duty shares?}
\author{}\date{}
\begin{document}\maketitle

\paragraph{Decision, observation, and installed-law alternatives.}
At one modeled cut, a known service through endpoint $A$ placed $q$ bytes
on its retained receiver ledger. Endpoint $B$ names either the
same physical ledger or a different unused ledger. It is not enough
to compare public endpoint strings. An authorized probe of nonzero
size observes its actual matching terminal receipt or lack of one,
then an old report of size $y>0$ is due through $B$.
There is no inter-probe expiry within the cut. The two installed
bindings remain fixed within a history, and both physical paths have
the same declared loss-free RF condition for the positive comparison.
The locally complete receipt watch is part of the fixture, not a
deduction from observing a remote silence.

\paragraph{One transmitter.}
With one initially uncharged sender and receive test ceiling $c$, a
separating probe must pass when $B$ is disjoint and fail when $B$
shares $A$. Refusal still spends its sender byte credit, because the
pinned native receipt admission occurs after the sender debit. The
minimum separating payload is $c+1-q$ for $0<q\le c$.
The old report must fit both the original shared receiver remainder
and the transmitter remainder after the probe. This is exactly
\[
             0<q,\qquad q+y\le c,\qquad y<q.
\]
For an adaptive sequence, let $a$ be the total of previous accepted
probes at its first differing observation, and $x>0$ the
differentiating payload. Before that event the policy observes the
same outcomes in both installations. Its transmit ledger has spent
$a+x$ at divergence. A distinguishing event needs $c<q+a+x$;
carrying the duty also needs $a+x+y\le c$ and $q+a+y\le c$.
The same lower bounds follow; the one-probe policy attains them using
$a=0$. No randomized guarantee over all allowed outcomes can avoid the
same constraint on each realized successful decision path. Lean checks
the exact finite arithmetic, and direct adaptive enumeration ties.

\paragraph{Different existing transmitters.}
A distinct probe sender of ceiling $p$ does not consume the old
sender budget. In the selected pinned implementation the receive
ledger is still that of the destination, but the ceiling for each
test is the current sender's. With an authorized distinct probe
sender and no change to the old sender's ceiling $c$, its minimum
separating load becomes $\max(1,p+1-q)$ when $q,p>0$.
Both services fit if this load plus $y$ is at most $c$ in the disjoint
receiver, and $q+y\le c$ in the shared receiver.
Merely separating the transmitters without changing their thresholds
can still waste the separate world's receive reserve.

For $(c,q,y)=(5,2,2)$, the same-sender four-byte probe distinguishes
but neither old duty completes. A separate equal-ceiling sender
still costs too much at the disjoint receiver, where the large probe
is accepted. A pre-existing authorized sender with ceiling two
separates with one byte, leaving both old deliveries possible.
These are three different available-action sets. The latter gain
is not against a direct scheduler already permitted to use the same
lower-ceiling sender. Nothing is remotely reconfigured in the fixture.

\paragraph{Native replay boundary.}
The replay executes the selected pinned SubT broker, radio and local
callback bodies on shared team radio states. An initial receipt makes
the retained byte ledger, the exact next callback is the probe
observation, the subsequent callback is the old service outcome,
and both come from the same executed version. It includes
18,049 native rows and 1,632 direct adaptive checks. An authored
RF-drop control applied only to the disjoint probe makes its observed
absence agree with the shared case. The positive comparison then
does not apply; neither a configured route nor a missing report alone
qualifies the binding. Early expiry, different observable path laws,
false receipt, or spare noninterfering resources define new problems.
This is a limited active-evidence development (R3), not an independently
selected mission, stock live communications experiment, or I1/I2
closure.
\end{document}
